\chapter{Arquitetura da BMS 4S}

\section{Visão geral do sistema}

A arquitetura proposta divide a solução em três camadas principais:

\begin{itemize}
  \item camada de medição e proteção primária;
  \item camada de controle e supervisão;
  \item camada de potência e comutação.
\end{itemize}

\section{Arquitetura de hardware}

Os blocos de hardware propostos são:

\begin{itemize}
  \item pack de células 4S;
  \item AFE dedicado;
  \item resistor shunt para medição de corrente;
  \item sensores NTC;
  \item MOSFET de carga;
  \item MOSFET de descarga;
  \item resistores de balanceamento;
  \item fusível.
\end{itemize}

\section{Papel do AFE}

O AFE ocupa posição central na arquitetura, pois fornece medição multicélula, suporte à proteção
e interface segura com o microcontrolador. Para uma arquitetura 4S, famílias como a BQ76920
apresentam bom alinhamento com os requisitos de medição e proteção \cite{ti_bq76920}.

\section{Papel do ESP32-S3}

O ESP32-S3 foi definido como controlador de alto nível, responsável por:

\begin{itemize}
  \item ler a telemetria consolidada;
  \item executar a máquina de estados;
  \item tratar falhas;
  \item controlar os caminhos de potência;
  \item expor interfaces de monitoramento e configuração.
\end{itemize}

\section{Estágio de potência}

O estágio de potência prevê caminhos distintos para carga e descarga. Essa separação facilita
o bloqueio seletivo de eventos de falha conforme o tipo de condição detectada.

\section{Sensoriamento de corrente e temperatura}

O sensoriamento de corrente é realizado por resistor shunt, enquanto a monitoração térmica
utiliza, no mínimo, dois termistores NTC: um próximo às células e outro próximo ao estágio
de potência.
